\begin{abstract}
Automating the digitization of Piping and Instrumentation Diagrams
(P\&IDs) into structured process graphs would unlock significant value
in plant operations, yet progress is bottlenecked by a fundamental data
problem: engineering drawings are proprietary, and the entire community
shares a single public benchmark of just 12 annotated images.
Prior attempts at synthetic augmentation have fallen short because
template-based generators scatter symbols at random, producing graphs
that bear little resemblance to real process plants --- and accordingly
yield only ${\sim}33\%$ edge detection accuracy under synth-only training.
We argue the failure is structural rather than visual, and address it
by introducing \ours{}, a corpus of 665 synthetic P\&IDs whose pipe
topology is seeded directly from real drawings.
Paired with a patch-based Relationformer adapted for high-resolution
diagrams, a model trained on \ours{} alone achieves
\textbf{63.8\,$\pm$\,3.1\% edge mAP} on PID2Graph~OPEN100 without
seeing a single real P\&ID during training, closing within 8~pp of
the real-data oracle.
These gains hold up under a controlled comparison against the
template-based regime of~\cite{sturmer2025}, confirming that generation
quality drives performance rather than model choice.
A scaling study reveals that gains flatten beyond roughly 400 synthetic
images, pointing to seed diversity as the binding constraint.
\ours{}, the generation code, and trained weights will be released
upon acceptance.
\end{abstract}
